\section{Semantics of Linear Programs}\label{sec:sem}

Next, we introduce the categorical semantics of parametric linear
programs (PLP). This section is self-contained. All results from the ordinary theory of linear programs, along with any categorical notions that are used here or later, are defined and proved within it.

Some results presented here, such as the existence of a standard form for PLP and certain
equivalences between their epigraph and value function, will be of fundamental importance
later when demonstrating the completeness of our algebra. Although those are know results from the literature, for the sake of comprehensibility we state and prove them over here.

Throughout this section, in order to facilitate the exposition of some ideas, we will often denote by $\llbracket = \rrbracket$ that, two objects are equal under the functor $\llbracket - \rrbracket$, i.e., a diagram $c \in Ob(FreeP_{(\Sigma,E)})$ is, under the interpretation map $c \llbracket = \rrbracket \llbracket c \rrbracket$.

\subsection{Preliminary Results in Linear Programming}

We begin with several classical results that establish the tight
relationship between piecewise affine convex functions, polyhedral
functions, and parametric linear programs. These will serve as the
semantic foundation for the graphical algebra introduced in the following
section.

\begin{definition}[Piecewise Affine Convex Function]
	A function $f : \mathbb{R}^n \to \mathbb{R}$ is called
	\emph{piecewise affine convex} if there exist finitely many affine
	functions $a_i^\top x + b_i$, $i = 1,\dots,k$, such that
	\[
		f(x) = \max_{i=1,\dots,k} \{ a_i^\top x + b_i \}.
	\]
\end{definition}

\begin{definition}[Polyhedral Function]
	A function $f : \mathbb{R}^n \to \mathbb{R} \cup \{+\infty\}$ is
	\emph{polyhedral} if its epigraph
	$\mathrm{epi}(f) := \{ (x,t) \mid t \ge f(x) \}$
	is a polyhedron.
\end{definition}

\begin{definition}[RHS Parametric Linear Program]
	A \emph{right-hand side parametric linear program} has value function
	$v(b) := \inf_{x} \{ c^\top x \mid Ax \ge b \}$.
\end{definition}

The following two propositions establish that all three notions above
are equivalent descriptions of the same class of functions. This
equivalence is the semantic cornerstone of the section: it is what
allows us to treat polyhedral functions and parametric linear programs
as morphisms in the same category. These results are classical (see
\cite{boyd2004convex, schrijver1986theory}); we include their proofs in
Appendix~\ref{app:proofs} for completeness.

\begin{proposition}\label{prop:pwa-polyhedral}
	A function is piecewise affine convex if and only if it is a polyhedral
	function.
\end{proposition}

\begin{proposition}\label{prop:pwa-value}
	A function $v : \mathbb{R}^m \to \mathbb{R} \cup \{+\infty\}$ is
	piecewise affine convex if and only if it is the value function of a
	right-hand side parametric linear program.
\end{proposition}

The following proposition is the key technical lemma for the completeness
proof of $\mathbf{GLP}$. It states that two parametric linear programs
with the same value function must share the same epigraph polyhedron,
providing the bridge between semantic equality and polyhedral equality
that the completeness argument requires.

\begin{proposition}\label{prop:equalepi}
	Let
	\[
		v_i(y) := \inf\{\mathbf{1}^\top x \mid A_i x \ge B_i y\},
		\qquad i=1,2,
	\]
	where $A_i \in \mathbb{R}^{m_i \times n_i}$ and
	$B_i \in \mathbb{R}^{m_i \times k}$. If $v_1(y) = v_2(y)$ for all
	$y \in \mathbb{R}^k$, then both programs admit a common representation
	through the same epigraph polyhedron:
	\[
		v_i(y) = \inf\{t \mid (y,t) \in E\},
		\qquad
		E := \mathrm{epi}(v_1) = \mathrm{epi}(v_2).
	\]
\end{proposition}

\subsection{The semantic category \texorpdfstring{$\mathbf{LPRel}$}{LPRel} and \texorpdfstring{$\mathbf{PolyFun}$}{PolyFun}}

With the semantic equivalences established, we now define the two equivalent semantic categories for this algebra. The first, $\mathbf{LPRel}$, describes morphisms as parametric linear programs; the second, $\mathbf{PolyFun}$, describes them as polyhedral functions. The equivalence between the two follows immediately from the propositions above. We refer only to $\mathbf{LPRel}$, but all results hold equally for $\mathbf{PolyFun}$.

The broadest semantic setting is the category of bifunctions, which is exactly the Cost/Reward instance of $\mathbf{Rel}(Q)$ for the Extended-Lawvere quantale. Stein and collaborators \cite{stein2024compositional, stein2025compositional} construct this category as the home of compositional convex optimization, where infimal convolution, conjugation, and variable elimination arise as categorical constructions; we recall only the bare construction we need.

\begin{definition}[Bifunction Category]
	The category $\mathbf{BiFunc}$ has sets as objects, and a morphism $F : A \to B$ is a function
	$F : A \times B \to \overline{\mathbb{R}}$.
	Composition is infimal convolution
	\[
		(G \circ F)(a,c) := \inf_{b \in B} \big( F(a,b) + G(b,c) \big),
	\]
	with identity $1_A(a,a') = 0$ if $a=a'$ and $+\infty$ otherwise ---
	precisely the characteristic function of the diagonal binary relation
	of Definition~\ref{def:binaryrel}. The tensor product is the Cartesian product on objects and addition on morphisms, with the singleton as unit.
\end{definition}

\begin{remark}
	The category $\mathbf{BiFunc}$ is exactly $\mathbf{Rel}(Q)$ for the Extended-Lawvere quantale
	$(\overline{\mathbb{R}}, \ge, +, 0)$: a $Q$-valued relation $R : A \times B \to \overline{\mathbb{R}}$ is precisely a bifunction, and the composition law $\bigvee_b R(a,b)\otimes S(b,c) = \inf_b (R(a,b)+S(b,c))$ is exactly infimal convolution. The Boolean-relations category embeds into $\mathbf{BiFunc}$ as a full subcategory via the indicator (characteristic) functor, so constraints and objectives are not fundamentally different objects: they are relations valued in different quantales, unified by the same compositional structure.
\end{remark}

Composition is therefore precisely variable elimination: composing two relations marginalises over the intermediate variable by minimising the combined cost\footnote{this notion is related to the idea of conditioning and marginalization in PPL~\cite{stein2024gqa, vandemeent2021introductionprobabilisticprogramming}}
\[
	(S\circ R)(x,z)=\inf_y \{R(x,y)+S(y,z)\}.
\]
Optimization is not imposed on the categorical structure from outside --- it emerges from composition itself. Restricting to proper, convex, lower semicontinuous bifunctions yields the convex subcategory $\mathbf{CxBiFn} \subseteq \mathbf{BiFunc}$~\cite{rockafellar1970convex}, which is closed under infimal composition and is the natural home for the optimization problems we study. Linear programs are a finitely-presentable class of morphisms inside it.

\begin{definition}[The prop $\mathbf{LPRel}$ of Linear Programs]
	The prop $\mathbf{LPRel}$ is defined as follows:
	\begin{itemize}
		\item Objects are natural numbers $n \in \mathbb{N}$, interpreted
		      as the space $\mathbb{R}^n$.
		\item A morphism $n \to m$ is a right-hand side parametric linear
		      program with value function
		      \[
			      (x,y) \mapsto \inf_z \{ c^\top z \mid Az \ge B[x\ y]^\top \},
		      \]
		      where $x \in \mathbb{R}^n$ and $y \in \mathbb{R}^m$.
		\item Composition is given by minimisation over the shared variable.
		\item The monoidal product on objects is addition $n \otimes m := n+m$,
		      and on morphisms is direct sum.
	\end{itemize}
\end{definition}

\begin{definition}[The prop $\mathbf{PolyFun}$ of polyhedral functions]
	The prop $\mathbf{PolyFun}$ is defined as follows:
	\begin{itemize}
		\item Objects are natural numbers $n \in \mathbb{N}$, interpreted
		      as the space $\mathbb{R}^n$.
		\item A morphism $n \to m$ is a polyhedral function
		      $f : \mathbb{R}^{n+m} \to \overline{\mathbb{R}}$.
		\item Composition is given by infimal convolution over the shared
		      variable.
		\item The monoidal product on objects is addition $n \otimes m := n+m$,
		      and on morphisms is $f \otimes g := f + g$.
	\end{itemize}
\end{definition}

\begin{remark}
	The two categories are isomorphic as PROPs: $\mathbf{PolyFun} \cong
		\mathbf{LPRel}$. This is a direct consequence of the equivalence
	between polyhedral functions and right-hand side parametric linear
	programs established in the propositions above. We will therefore
	use the two names interchangeably, choosing whichever description is
	more natural in a given context.
\end{remark}

